\chapter{Scalar arithmetic}\label{ch:scalar-arithmetic}

Beside the tensor unit, every G17 lane runs an ordinary scalar ALU: integer arithmetic on 32-bit
registers and 16-bit halves, fp32 and fp16 arithmetic, conversions, compares that write flags or select values, and five
single-instruction special functions. A compiler targeting this machine must reproduce the ALU bit for bit, because a
tensor kernel's address arithmetic, quantization and epilogues all run here.

Apple does not document this ALU. Many of its forms were first named from Apple's decoder
tables and compiler output, and a later audit of those names against execution corrected 7 of 100 uniquely fitted forms
(MM~\mm{25.90}). Each semantic below was measured on inputs chosen so that the rival readings give different answers.
Several results bear directly on code generation. The integer condition codes are a structured field, and one code
that every source called "equal" is a bit test (\cref{sec:compares-flags-selects}). The machine has three subnormal
policies, one per execution path, so a value can change when it passes from one path to the next; and both fused
multiply-adds, fp32 and fp16, round once, proven on inputs where once and twice differ
(\cref{sec:fp32-arithmetic,sec:fp16-arithmetic}). On positive normal inputs the hardware reciprocal
square root is an exact, tabulable function with a scaling law, while hardware exp2 is deterministic but has no compact
description; the difference determines how kernels that use them are checked (\cref{sec:special-function-instructions}).

The full per-opcode table, with lengths and standing, is \cref{app:b}.

Conventions:

\begin{itemize}
\item Opcodes are cited as glossary name then number, "\op{10282}"; "/N" is the encoded length in bytes, and a
  semantic belongs to the (opcode, length) forms that ran. Forms Apple's compiler emits that have no glossary entry are
  given a label and marked "(no glossary entry)".
\item Rn is general register n (R0–R125 usable, \cref{ch:register-file}); RnL and RnH are its 16-bit halves. a32 and a16 are 32- and
  16-bit sources; zext16 and sext16 are zero and sign extension of a half; K8 is an 8-bit float-coded immediate
  (\code{isa/g17-float-immediate.toml}).
\item Each register source carries a \textbf{modifier operand} (\cref{sec:operand-modifiers}): negate (2) and absolute value (4) act on values and
  are used here; release (16) and keep (32) govern lifetime (MM~\mm{6}).
\item FTZ: subnormal inputs and results flushed to zero. RNE: round to nearest, ties to even. A canonical NaN is the single
  pattern a form returns for every NaN input.
\end{itemize}

Many forms below can release a source. Treat a later read of a released register as "zero or the
old value", and read one only where both outcomes agree (MM~\mm{25.95}, MM~\mm{25.145}, MM~\mm{25.145.4}, \code{agxforge/g17/releasecheck.py};
the rule and its evidence are \cref{sec:release-does}).

\section{Integer arithmetic}\label{sec:integer-arithmetic}

Integer results are the low 32 bits of the exact result unless a form says otherwise. The families are conventional.
Their edge cases below were each measured on an input where the candidate readings disagree (\cref{app:b}
lists every form and its source).

\subsection*{Add, subtract and saturating forms}

\op{10282}, \op{10279} with an 8-bit immediate,
\op{11667} and \op{11666} wrap. The forms that mix widths zero-extend, and
each was separated from sign extension by an input with the half's top bit set:

\begin{itemize}
\item \op{10283}: d = a32 + zext16(b) (MM~\mm{25.145.1});
\item \op{10285}: d32 = zext16(a) + b32. Its census operands were special-register values below $2^{15}$,
  where zero and sign extension agree, so a lane kernel with a full-width operand was needed (MM~\mm{25.195});
\item \op{10286} widens: 0xFFFF + 1 = 65536, refuting a 16-bit wrap;
\item \op{10289} does wrap at 16 bits, d16 = a16 + imm (4,096 lanes; "immediate + 1" refuted,
  MM~\mm{25.195}).
\end{itemize}

\op{10239} and \op{11624} clamp; the receipt that fixed them chose unsigned
saturation over signed saturation and wrapping on 32 lanes (MM~\mm{25.145}). Their signedness is not an opcode: in the corpus,
the saturating add carries source modifier 16 (32 for a high-half source) in every unsigned kernel and 24 (40) in every signed one
(MM~\mm{25.59}), and \op{16805} executed only with modifier 24 (MM~\mm{25.145.1}). Reading
the value 8 as a signedness bit fits both; that it is a general integer-modifier rule is \emph{open}.

\subsection*{Multiply}

\op{10825} and \op{10826} keep the low 32 bits. Around them sits a
family of immediate and half-width variants (\cref{app:b}) whose decoder names were repeatedly wrong: the three subtracting
forms (\op{11059}, \op{11060}, \op{11063}) carry the map name \code{madd}, and three 64-bit product forms carry \code{mulhi} though they compute full products
(glossary). \op{10838} computes sext16(a) * zext16(b) + c32, 60 of 60 over seven
records (MM~\mm{25.59}). There is no integer divide: \code{A[t] / (B[t] | 1)} lowers to a reciprocal sequence through \op{3658} that touches fourteen forms (Apple's compiler emits, MM~\mm{25.80}).

\subsection*{Logic, shifts and bit operations}

\begin{itemize}
\item \op{14392} uses the amount's low 7 bits and gives 0 when the amount is 32 or more, not
  a shift modulo 32 (measured, glossary).
\item \op{9986} returns 0xFFFFFFFF for an input of 0 (the rival readings, 0 and a leading-zero
  count, were refuted); clz is derived from it (receipt fit, MM~\mm{25.145}).
\item The arithmetic shift right immediate sign-fills; the other right shifts, by register and by immediate, are logical.
\item The 4-byte register forms release their sources. \op{424}/4, \op{13575}/4 and \op{17771}/4 have no lifetime bit and always release both sources; all three matched Apple's compiler element for
  element on disjoint operands (measured, MM~\mm{24.8} row T4, MM~\mm{25.77}). A source needed again requires the /10 form or a
  copy.
\item \op{428} has a uniform mask. Authored with a register as its mask, it returned zero on
  every lane, because its second operand is a uniform halfword of the kernel's constant pool and not a
  register; the pool sits after the binding pointers (uniform K is pool halfword K - 4 × buffers). Zeroing or
  swapping that pool fails all 2,048 outputs of a q4 quantized matrix-vector chain (whole-program, MM~\mm{25.139.1},
  MM~\mm{25.141.17}, MM~\mm{25.145}). The pool format is \cref{ch:kernel-metadata-image}'s. Extracting a half from a register takes \op{426} instead, which reads one half with its source kept and is exact for L and H (MM~\mm{25.139.1}).
\end{itemize}

\subsection*{Moves and immediates}

\op{586} and \op{590} copy a register or one half.
\op{554} has run only with immediate 0 (1,228 of 1,231 corpus uses are 0), and \op{555} is claimed only for 0 (MM~\mm{25.145.1}, MM~\mm{25.184}). \op{11842}, 10 bytes, carries a full 32-bit immediate. \op{13574} reaches
R64 and above and can carry a load wait, so OR with 0 is the register copy into the upper file (measured,
glossary). Immediate size selects the form (Apple's compiler emits, every probe executed, \code{isa/g17-comparison-semantics.json}):
a select keeps both-immediate arms 0..255 in an immediate-arm select (\op{11373}, no glossary entry) and moves 256 and above,
and every negative value, to \op{11375}; a compare against an immediate holds 0..256
in \op{11365} and moves 257 and above to the four-register form.

\section{Compares, flags and selects}\label{sec:compares-flags-selects}

A comparison and its select are one instruction (Apple's compiler emits). \code{a < b ? x : y} on int emits a single \op{11372}/14 or a four-register select, and no separate compare
(\code{isa/g17-comparison-semantics.json}). Compares that must survive as predicates write a flag register: \op{10369}, "the flag compare" below, is the loop-exit compare, /10 for a \code{uint} counter and /6 for an \code{int} one (MM~\mm{25.73}).
The flag registers FLAG0-FLAG14 are an allocatable class (\cref{ch:register-file}).

\subsection*{Integer condition field}\label{sub:integer-condition-field}

The decoder prints an integer condition as a number from 8 to 15. Authoring \op{11313} at
each printed code and reading the flag on four operand pairs fixed the decode (measured, \code{isa/g17-condition-codes.toml}):

\begin{lstlisting}
printed cc = 8 | (signed << 2) | (gt << 1) | lt        equality: lt = gt = 0
\end{lstlisting}

\widetable
\begin{xltabular}{\linewidth}{@{}lXllll@{}}
\caption{Integer condition codes by printed number, with the result of 32-bit compare, OR-chained into a flag on four operand pairs.}\label{tab:integer-condition-codes}\\
\toprule
printed & relation & (10,20) & (20,10) & (10,10) & (−5,3) \\
\midrule
\endfirsthead
\tablecontinued{6}
\toprule
printed & relation & (10,20) & (20,10) & (10,10) & (−5,3) \\
\midrule
\endhead
8 & unsigned equal & false & false & true & false \\*
9 & unsigned less & true & false & false & false \\*
10 & unsigned greater & false & true & false & true \\*
12 & signed equal & false & false & true & false \\*
13 & signed less & true & false & false & true \\*
14 & signed greater & false & true & false & false \\*
15 & predicted "not equal"; measured otherwise & false & false & true & true \\*
\bottomrule
\end{xltabular}

In \cref{tab:integer-condition-codes}, signed and unsigned order disagree on the pair (−5, 3), so that column identifies bit 2 as signedness rather than a
correlated bit. The decode predicted that codes 11 and 15 mean not-equal (lt or gt), but code 15 is true on
(10,10) and (−5,3) only, which is the bit test described below.

An encoder writes a field code, not the printed number. For the immediate-arm compare-select the field
codes are 0 equal, 1 unsigned less-than and 5 signed less-than (printed = 8 \textbar{} field), and field 4 behaves as 0, as the
decode predicts. Every code returns exactly 0 or 1, so the remaining seven relations derive by swapping operands and by
complementing with one XOR: ten relations at four orderings, 40 of 40 (measured, \code{ledger/notes.toml}
{[}g17-all-ten-relations-from-three-codes{]}).

\subsection*{Select field code 1}\label{sub:field-code-1}

\op{11375} computes d = (a cc b) ? x : y with two field codes. Code 0 is a > b. Code
1 was recorded as a == b, from witness pairs (1,2), (2,1), (1,1) and (5,5), and shipped in this compiler as
\code{CSEL\_CC['eq']}. Those four pairs cannot tell equality from a bit test: the unequal pairs share no bit and the equal pairs
are nonzero, so both readings give the same answers. The error was found when a switch, lowered to a chain of selects that
tested a 0/1 compare value against 0, chose the wrong arm on every lane. The discriminating run (measured, MM~\mm{25.90};
\code{ledger/claims.toml} {[}g17-csel-code-one-is-a-bit-test{]}): x = 0..31 on 32 lanes,

\begin{itemize}
\item against 16: selected exactly x ≥ 16, in both operand orders (symmetric, so not an order relation);
\item against 0: never selected (equality would select lane 0);
\item against 5: selected exactly the 24 lanes with x \& 5 != 0 (equality gives one lane, ≥ gives 27).
\end{itemize}

Field code 1 therefore computes (a \& b) != 0. \emph{Correction:} the \code{a == b} reading (MM~\mm{25.90}'s "Was" column) is wrong; cc now names
code 1 \code{test} and refuses a select that asks for \code{eq}.

The four-register select's field codes 0 and 1 print as 14 and 15; \code{tools/g17emu.py} uses
the printed numbers and the ledger the field numbers. The printed-15 truth row of \cref{tab:integer-condition-codes} (false, false, true,
true) is exactly the bit test: 10 \& 20 = 0, 20 \& 10 = 0, 10 \& 10 = 10, and −5 \& 3 = 3. g17emu reads printed 14 as a > b
and refuses inputs at or above $2^{31}$, matching the ledger's open signedness for code 0. Field 1 (printed 15) can be
relied on as the bit test, never equality. \emph{Open:} the bit test on operands with high bits set (the
test used x < 32), and code 0's signedness above $2^{31}$. \op{11507} is plain
equality; on separating cases it fits a == b alone (MM~\mm{25.90}).

\subsection*{Flags, chains and immediates}\label{sub:flags-chains-immediates}

\begin{itemize}
\item \code{\&\&} and \code{||} chain through flag-reading compares: \op{11310}
  flag = (x cc K) ? c16 : 0, and the OR-chained compare flag = (x cc y) ? T : flag\_in (measured, glossary; more chain forms in \cref{app:b}).
\item \code{if (x == 0)} and \code{if (x != 0)} compile to the same flag compare with printed code 12 and differ only in the exec form's
  invert bit (decoder, \code{isa/g17-exec-mask.txt}, \cref{sec:execution-mask-family}).
\item Apple's compiler never emits ≥ or ≤ against an immediate: \code{x >= K} becomes \code{x > K-1} and \code{x <= K} becomes \code{x < K+1}
  (\code{isa/g17-scalar-isa.toml}).
\item A flag's numeric value is form-specific. Read back through \op{14099} after one compare form
  (\op{10379}, no glossary entry), a flag read 1 for true and 2 for false (\code{ledger/notes.toml}); the authored OR-chained
  compare read 2 for true and 1 for false (\code{isa/g17-condition-codes.toml}). Both are measured on their own forms. Use a flag only
  as a predicate.
\item The flag compare's /10 form has a second operand layout selected by operand 1's $2^{31}$ bit (1,421 vendor instances against 19,260); no
  source shape reaches it and its meaning is \emph{open} (MM~\mm{25.73}, MM~\mm{25.76}).
\end{itemize}

\subsection*{Floating-point compares}\label{sub:floating-point-compares}

Every ordered relation is false when either operand is NaN, != is true on NaN, and −0 equals +0, for float, half and
bfloat (measured, \code{isa/g17-comparison-semantics.json}). \op{9700} is the four-register float
select; its codes are 0 ==, 1 \textless, 2 \textgreater, 4 !=, 5 ≥, 6 ≤, and Apple's compiler writes 3 for \code{min(float, float)} and 7 for
\code{max} (\code{agxforge/g17/cc.py}). For 3 and 7 the hardware tests a first: a NaN a selects y, then a NaN b selects x, and with
both NaN it returns y. A model that tested b first disagreed with the receipts on lanes 7 and 23 (measured, MM~\mm{25.145}). \code{llvm.maxnum(x, 0)} returns 0 for NaN x, so relu(NaN) = 0 (MM~\mm{6}).

\section{fp32 arithmetic}\label{sec:fp32-arithmetic}

\op{998}, \op{3290} and \op{2190} are RNE. There is no subtract instruction: \code{a - b} is the add with the negate modifier on b, exact since negation is exact
(Apple's compile, MM~\mm{25.180}). Literal operands select different forms (\op{2192}; \op{2198}; \op{2200}; \op{1006}), and \code{a * 3.5f} selects a form distinct from \code{a * b} (MM~\mm{25.73}).

Two fp32 FMA forms were tested for single rounding, on inputs where fused and unfused evaluation round differently.

\begin{itemize}
\item The fp32 fused multiply-add: 0 of 16 rounding-witness lanes double-rounded (MM~\mm{25.90}). Its four-byte form needs the
  accumulator and destination in one register.
\item The literal-addend FMA, from Apple's \code{fma(a[i], b[i], 1.0f)} over 4,096 lanes whose draws put the +1 inside
  the product's discarded bits: on the 3,186 lanes where once and twice differ, the fused model matches all and the
  unfused model misses all (measured, MM~\mm{25.184}, \code{isa/g17-rounding-results.json}).
\end{itemize}

Subnormals are handled three ways. The fp32 ALU flushes subnormal inputs and results to zero (fadd, fmul, the immediate fma
forms, the reciprocal); the tensor unit's fp16, bf16 and fp32-accumulator paths flush nothing; and the measured half ALU
forms keep half subnormals (\cref{sec:fp16-arithmetic}). An MMA can produce a subnormal that the next fp32 ALU
instruction to read it flushes to zero (measured, MM~\mm{5}, MM~\mm{17} rule 22). The flush keeps the sign. A GELU epilogue that passed within its
bound at M16 to M64 failed in 57 of 65,536 elements at M512, all with x near −52, where 1/(1 + 2\^{}t) is about 3.6e-39;
the GPU returned −0, and with FTZ in the reference every arm passed (measured, MM~\mm{25.109}).

fp32 NaN results are the canonical 0x7fc00000 (g17emu EVIDENCE; MM~\mm{5}). \op{904}
is encoded as sat(a + K) with K = ±0 and maps −0 to +0, values above 1 to 1 and negatives to 0 (32 lanes); NaN into it is
unmeasured. \emph{Correction:} the map's "fsat" is \op{1062}, a cross-lane form (\cref{sec:other-cross-lane-forms});
saturate is the form above (MM~\mm{25.90}). \op{3818} keeps the sign of a zero result (−0.3 → −0), and \op{3770} is Metal \code{rint}.

\code{air.dot.vNf32} is an fp32 multiply of lane 0 then one fma per lane, left to right; a
pairwise order differs from Apple's on 66 of 256 random lanes, while an fp64 model of Apple's chain agrees on all (MM~\mm{25.180}).
Under default fast math \code{x != x} folds to false, so a NaN guard must be \code{isnan(x)} or
\code{(as\_type<uint>(x) \& 0x7fffffff) > 0x7f800000} (measured, MM~\mm{25.50}).

\section{fp16 arithmetic}\label{sec:fp16-arithmetic}

The measured half forms run on 16-bit halves of the general registers, keep subnormals, and return the half canonical
NaN 0x7e00.

\op{798} (\code{ffma.f16.l12}, 12 bytes) computes binary16 a * b + c rounded once (measured,
MM~\mm{25.196}).

\begin{itemize}
\item Operands: destination 0; sources 2, 4 and 6; their lifetimes 3, 5 and 7. A source is a low half (register id 425 + n)
  or a high half (281 + n) of word register n; Apple's own \code{fma(half, half, half)} reads a high half (r283). Authoring
  Apple's operands reproduces Apple's bytes, \hash{3800060a2320a40245018000}.
\item The test: 196,608 triples in three cases (the whole finite binary16 range by bit pattern, and two cancelling cases with
  24,000 subnormal results each). 0 differ from binary16(a * b + c) computed exactly; the multiply-then-add control
  differs on 6,246 to 54,951 per case; with both addends taken from the two halves of one packed word, 0 of 196,608
  differ.
\item Encoding holes (decoder): the low-half destination cannot author r18, r19, r26, r27, r50, r51, r58 or r59, operand 4
  cannot author r17L and operand 6 cannot author r18L. The load-wait field is not located.
\end{itemize}

\textbf{Worked example: exact 4-bit dequantization in two fp16 fmas} (\emph{this compiler}; executed bit-exact on hardware,
MM~\mm{25.196}). Integer code packs a 4-bit weight q into a half's mantissa under the bits of 1024.0h, so the half's value is
h = 1024 + 2\^{}k q (k = 0, 4 or 6, masks 0x000F00F0 and 0x03C000F0). One fp16 fma recovers q exactly as h * 2\^{}-k - $2^{10-k}$:
every intermediate is a small integer, so the single rounding is exact. A second fp16 fma, fma16(q, s, b), applies the
scale and bias with one rounding and writes the weight straight into the half A fragment of a tensor tile. On a sample
of 29,360,128 real weights the result equals the old fp32-then-fp16 dequant on every weight (the fp32 sum of a bf16
product and bias is already exact), so the model's logits do not move; the projection gets 10 to 12 percent faster because
the dequant leaves the fp32 pipe (128 fp16 fmas per trip, none of the integer-to-float, fp32 fma or widening forms).

The mixed-width forms below are all measured.

\begin{itemize}
\item \op{776}: half(a + K8). A NaN input returns 0x7e00, not its payload (104 and 90 NaN lanes);
  subnormals are kept (a flushing model misses 965 lanes of an \code{h + 1/64} kernel built to reach subnormal results);
  rounding once and through fp32 are the same function on all 65,536 halves and every positive immediate, checked
  exhaustively on the CPU, so this form needs no rounding test (MM~\mm{25.195}).
\item \op{1014}: half(a32 + b32), one rounding, overflow to inf, NaN → 0x7e00, a tiny negative → −0.
  \op{1015} keeps half subnormals. \op{3306}: a width receipt's
  0x3e00 rules out a bf16 result (MM~\mm{25.145.1}).
\item Apple's compiler selects the fp16 fma (/12) for \code{fma} on half and a different form (\op{862}/10, no glossary entry) for \code{a * b}
  (MM~\mm{25.73}).
\end{itemize}

\section{Conversions}\label{sec:conversions}

\begin{xltabular}{\linewidth}{@{}>{\hsize=0.69\hsize}XX>{\hsize=1.41\hsize}X>{\hsize=0.86\hsize}X@{}}
\caption{Rounding and edge behaviour of the measured conversions.}\label{tab:conversions}\\
\toprule
Conversion & Form & Rounding and edges & Standing \\
\midrule
\endfirsthead
\tablecontinued{4}
\toprule
Conversion & Form & Rounding and edges & Standing \\
\midrule
\endhead
fp32 → half & \op{1016} & IEEE RNE; zero mismatches over 1,510,400 fp32 patterns, subnormal
results included; NaN → 0x7e00 & measured (MM~\mm{6}) \\*
fp32 → bfloat & \op{1048} is the candidate & RNE except that an fp32 subnormal input returns a zero with its sign
(5,879 of 5,912 subnormal patterns); NaN → 0x7fc0;
denormal enable and fast-math off change nothing & function measured (MM~\mm{6}, MM~\mm{17} rule 21); instruction identity
unresolved \\*
half → fp32 & \op{1004} & exact & measured \\*
int32/uint32 → fp32 & \op{11179}/10 & RNE; operand 2: 4 unsigned, 5 signed & measured against 8 rivals (MM~\mm{25.90}) \\*
int16/uint16 → fp32 & \op{11180} & exact; code 2 unsigned, 3 signed & measured (MM~\mm{25.195}) \\*
uint32 → half & \op{11182} & RNE; past half's range +inf, not 65504 (2,992
lanes) & measured (MM~\mm{25.184}) \\*
fp32 → int32/uint32 & \op{9320}/10 & mode operand 1: truncate toward zero, saturate, NaN → 0 & measured (MM~\mm{25.186}) \\*
fp32/half/bf16 → fp8 & \op{13618} & RNE, no saturation, every NaN → canonical NaN, over all
$2^{32}$ fp32 patterns & measured \\*
\bottomrule
\end{xltabular}

Remaining conversions (16-bit integer to half, half to integer, uniform-destination forms, fp8 unpack) are in \cref{app:b}.

Apple's float-to-int casts \code{(uint)x} and \code{(int)x} are \op{9320} with code 4 or 5 and mode
operand 1; the instruction bytes are \hash{2f00001a2200ae02b003} and \hash{2f00001a2200ae027001}. Over 4,096 lanes each, drawn to
cover fractions of both signs, ties, negatives, values past $2^{31}$ and $2^{32}$, NaN, +-inf and ±0: for \code{(uint)x},
truncate-and-saturate matches every lane, RNE-and-saturate misses 326 and truncate-and-wrap misses 1,974; for \code{(int)x},
truncate-and-saturate matches every lane and RNE-and-saturate misses 653 (measured, MM~\mm{25.186}). On \code{(uint)(-5.5)}, which
Metal leaves undefined, the instruction gives 0; this compiler now emits it in place of its own earlier decomposition
(\cref{app:a}).

Because the conversion saturates, it maps NaN to 0, +Inf to 2,147,483,647 and −Inf to −2,147,483,648 on its own, and the
usual "safe" \code{maxnum}/\code{minnum} clamp placed before it changes the NaN result: it sends NaN to the clamp's lower bound (−127 for an
int8 quantizer) (measured, MM~\mm{25.50}, MM~\mm{17} rule 17).

Signedness codes are per form: 4/5 on the 32-bit integer-to-fp32 and fp32-to-integer conversions, 2/3 on the 16-bit
integer-to-fp32 conversion. On the integer-to-float forms operand 2 is signedness: an unsigned code converts a negative
int32 to a value near $2^{32}$ (MM~\mm{25.90}). Canonical NaNs by destination: fp32 0x7fc00000, half 0x7e00, bfloat 0x7fc0.

\section{Special-function instructions}\label{sec:special-function-instructions}

\op{3850}, \op{1272} and \op{2570} are what Apple's compiler selects for \code{rsqrt},
\code{exp2} and \code{log2}, each at 10 bytes (MM~\mm{25.73}); \op{3658} and \op{3978} complete the
set. sqrt-safe rsqrt returns 1.0 at 0 where the seed returns +inf, so \code{fast\_sqrt}'s \code{x * rsqrt(x)} gives 0 at 0 where the
seed would give NaN (MM~\mm{6}). A \op{3946} is emitted by \code{cos}, \code{sin} and \code{cospi}.

Each special function was dispatched with a sentinel over every output, compared with the correctly
rounded value, and run twice (\cref{tab:special-functions}). The columns come from three different input populations, so percentages should not be
compared across columns.

\begin{xltabular}{\linewidth}{@{}>{\hsize=0.66\hsize}X>{\hsize=0.72\hsize}X>{\hsize=1.33\hsize}X>{\hsize=1.28\hsize}X@{}}
\caption{Accuracy and structure of the special-function instructions. Each column uses its own input population.}\label{tab:special-functions}\\
\toprule
Function & 1,024-input probe (MM~\mm{25.136.3}) & Dense or full measurement & Structure \\
\midrule
\endfirsthead
\tablecontinued{4}
\toprule
Function & 1,024-input probe (MM~\mm{25.136.3}) & Dense or full measurement & Structure \\
\midrule
\endhead
\op{3850} & 974 correctly rounded, 39 one ulp high, 11 low & all 3 × $2^{24}$ inputs of three exponent bands: correctly rounded on
16,051,395 of 16,777,216 (parity, mantissa) keys, +1 ulp on 628,246, −1
ulp on 97,575; at most 0.82 ulp (MM~\mm{25.141.16}) & exact function: seed(x * 4\^{}k) == seed(x) * 2\^{}-k bit for
bit, 16,777,216 of 16,777,216 at k = −8 and k = +5; \opn{3850}(0) = +inf \\*
\op{3658} & 931 exact, 91 high, 2 low &  & within one ulp; a subnormal result flushes to signed zero (MM
\mm{25.109}) \\*
\op{1272} & 632 exact, 387 high, 5 low & 496,131 inputs on t in {[}-127, 1{]}: 64.2 percent exact, 35.6 percent
+1 ulp, 0.17 percent −1 ulp; within 1.03e-7 relative over t in {[}-30,
0{]} (MM~\mm{25.144.2}) & deterministic (0 of 496,131 differ across dispatches); no
compact law \\*
\op{2570} &  &  & Apple's compiler emits it; accuracy not measured \\*
\bottomrule
\end{xltabular}

The rsqrt result came from sweeping whole bands of $2^{24}$ inputs at three
exponents (E = 127, 111 and 137) and testing the law seed(x * 4\^{}k) == seed(x) * 2\^{}-k across them. It held on every input
at both shifts, so one table over (exponent parity, mantissa) reproduces the instruction everywhere on positive normals; the table
is \code{isa/g17-rsqrt-seed.npz} and reproduces all 3 × $2^{24}$ hardware values. The same method does not work for hardware exp2. The
analogous law exp2hw(t + n) == exp2hw(t) * 2\^{}n holds on only 88 percent of inputs at n = 1 and 4 percent at n = 40,
because the fp32 add that forms t + n re-rounds the fraction and the instruction depends on its full input. The best constant-bias model,
f32(exp2(t) + 0.36 ulp), matches 97.1 percent, and a dense table would need about $10^{9}$ entries (MM~\mm{25.144.2}).

Kernels that use rsqrt and reciprocal are checked differently from kernels that use exp2.

\begin{itemize}
\item rsqrt and reciprocal are made correctly rounded on the GPU by an exact integer midpoint test. For a candidate a, the
  midpoint to its successor is (2 Ma + 1) $2^{Ea - 151}$; 1/sqrt(x) lies below it exactly when
  Mx (2 Ma + 1)\^{}2 > $2^{452 - Ex - 2 Ea}$, and 1/g exactly when Mg (2 Ma + 1) > $2^{301 - Eg - Ea}$. The products are formed in
  15-bit limbs with \op{10825}; a tie is impossible because 2 Ma + 1 is odd and at least $2^{24}$; and
  y0 - 1 + {[}f above mid(y0 - 1, y0){]} + {[}f above mid(y0, y0 + 1){]} is correctly rounded whenever the seed is within one ulp.
  On the probe the corrected values equal the correctly rounded ones on 1,024 of 1,024 (measured, MM~\mm{25.136}). Scope:
  positive normal inputs with a normal result; a seed two ulps off is not repaired. Alternatively, the seed alone is
  reproducible from its table.
\item exp2 has neither route, since no cheap exact test decides 2\^{}t against a midpoint, and no compact model reproduces hardware
  exp2. A kernel that uses hardware exp2 (attention softmax, GELU) is checked against an \emph{enclosure}: each output must
  lie inside the interval a reference derives when exp2 and reciprocal may each be −1, 0 or +1 ulp off (MM~\mm{25.109},
  MM~\mm{25.114}, MM~\mm{25.144.2}), and is never claimed bit-exact. Where bit-exactness is required this compiler uses \code{exp2\_soft},
  a fixed RNE sequence (clamp to {[}-125, 125{]}, split off rint(t), a degree-7 Taylor polynomial in Horner form, the integer
  part added to the exponent field), within one ulp of 2\^{}t and exact on 90.6 percent of 800,000 inputs (MM~\mm{25.136}). For
  replay the emulator reads hardware exp2 and hardware reciprocal from sparse tables of the inputs real programs feed them and refuses any
  other input (MM~\mm{25.145.4}).
\end{itemize}

Library functions have these measured edges (MM~\mm{6}): \code{air.fast\_tanh(t)} returns NaN for t ≥ 44.364; \code{air.fast\_divide} and \code{air.erf} do not
lower on this toolchain.

\section{Scalar dependences and issue}\label{sec:dependences-issue-scalar}

Which producers need a wait, and what an unwaited consumer reads, are \cref{sec:scoreboard}'s: ordinary ALU to ALU is interlocked
(52 of 52 adjacent pairs, MM~\mm{25.90}), late values (loads, atomic returns and the other memory producers) are not, and
whether an ALU consumer of an MMA result is interlocked is \emph{open} (MM~\mm{6}).

\begin{itemize}
\item Without the wait, scalar consumers read the old contents: a compare read the register's previous value
  (\code{isa/g17-scalar-isa.toml}), the six-byte fp32 fma got 0 of 32 lanes right (MM~\mm{25.90}), and an indirect load differing in
  one wait bit got 1 of 32 lanes right without it and 32 of 32 with it (MM~\mm{25.145}).
\item Some forms cannot carry the wait. On \op{11372} the load-wait position
  of the other forms is an opcode bit; on \op{798} the field is not located. \emph{This compiler} routes
  such an operand through one waited copy, \op{10279} with 0 (\code{agxforge/g17/cc.py}, MM~\mm{25.196}).
\item A load whose result is never read is still in flight. Reusing its destination while the load was outstanding let it land late and
  overwrite the new value (rows 2 and 9 of 16 failed); a register an outstanding load will fill must not be written
  (measured, MM~\mm{25.114}).
\item Special-register reads publish a slot but need no wait in the measured form (\cref{sec:read-latency-whether}).
\item Issue intervals, latency classes and the absence of fp32/int32 dual issue are \cref{ch:performance}'s.
\end{itemize}

\textbf{Safe assumptions.} The numerics above hold for the forms and lengths named, within each measurement's stated
domain. For a released register nothing holds beyond "zero or the old value".

\textbf{Not known.} NaN into saturate or truncate; negate and abs on the float select; the bit test with high-bit operands;
the mode-SET layout of \op{10369}; hardware log2's accuracy; the trig op's
function; the integer-modifier signedness bit beyond two opcodes; half results where rounding once and through fp32 differ
on forms other than \op{776} and \op{11182}.

\textbf{Evidence.} MM~\mm{5}, MM~\mm{6}, MM~\mm{17} (rules 17, 21, 22), MM~\mm{24.8} (row T4), MM~\mm{25.50}, MM~\mm{25.59}, MM~\mm{25.73}, MM~\mm{25.76}, MM~\mm{25.77},
MM~\mm{25.80}, MM~\mm{25.90}, MM~\mm{25.95}, MM~\mm{25.109}, MM~\mm{25.114}, MM~\mm{25.136}, MM~\mm{25.139.1}, MM~\mm{25.141.16}, MM~\mm{25.141.17}, MM~\mm{25.144.2},
MM~\mm{25.145}, MM~\mm{25.145.1}, MM~\mm{25.145.4}, MM~\mm{25.180}, MM~\mm{25.184}, MM~\mm{25.186}, MM~\mm{25.192}, MM~\mm{25.195}, MM~\mm{25.196};
\code{isa/g17-condition-codes.toml}, \code{isa/g17-comparison-semantics.json}, \code{isa/g17-scalar-isa.toml}, \code{ledger/notes.toml},
\code{ledger/claims.toml}, \code{tools/g17emu.py}, \code{agxforge/g17/cc.py}.
